\documentclass{amsart}
% For dealing with references we use the comment environment
\usepackage{verbatim}
\newenvironment{reference}{\comment}{\endcomment}
%\newenvironment{reference}{}{}
\newenvironment{slogan}{\comment}{\endcomment}
\newenvironment{history}{\comment}{\endcomment}

% For commutative diagrams we use Xy-pic
\usepackage[all]{xy}

% We use 2cell for 2-commutative diagrams.
\xyoption{2cell}
\UseAllTwocells

% We use multicol for the list of chapters between chapters
\usepackage{multicol}

% This is generally recommended for better output
\usepackage{lmodern}
\usepackage[T1]{fontenc}

% For cross-file-references

% Package for hypertext links:
\usepackage{hyperref}

% For any local file, say "hello.tex" you want to link to please

% Theorem environments.
%
\theoremstyle{plain}
\newtheorem{theorem}[subsection]{Theorem}
\newtheorem{proposition}[subsection]{Proposition}
\newtheorem{lemma}[subsection]{Lemma}

\theoremstyle{definition}
\newtheorem{definition}[subsection]{Definition}
\newtheorem{example}[subsection]{Example}
\newtheorem{exercise}[subsection]{Exercise}
\newtheorem{situation}[subsection]{Situation}

\theoremstyle{remark}
\newtheorem{remark}[subsection]{Remark}
\newtheorem{remarks}[subsection]{Remarks}

\numberwithin{equation}{subsection}

% Macros
%
\def\lim{\mathop{\mathrm{lim}}\nolimits}
\def\colim{\mathop{\mathrm{colim}}\nolimits}
\def\Spec{\mathop{\mathrm{Spec}}}
\def\Hom{\mathop{\mathrm{Hom}}\nolimits}
\def\Ext{\mathop{\mathrm{Ext}}\nolimits}
\def\SheafHom{\mathop{\mathcal{H}\!\mathit{om}}\nolimits}
\def\SheafExt{\mathop{\mathcal{E}\!\mathit{xt}}\nolimits}
\def\Sch{\mathit{Sch}}
\def\Mor{\mathop{\mathrm{Mor}}\nolimits}
\def\Ob{\mathop{\mathrm{Ob}}\nolimits}
\def\Sh{\mathop{\mathit{Sh}}\nolimits}
\def\NL{\mathop{N\!L}\nolimits}
\def\CH{\mathop{\mathrm{CH}}\nolimits}
\def\proetale{{pro\text{-}\acute{e}tale}}
\def\etale{{\acute{e}tale}}
\def\QCoh{\mathit{QCoh}}
\def\Ker{\mathop{\mathrm{Ker}}}
\def\Im{\mathop{\mathrm{Im}}}
\def\Coker{\mathop{\mathrm{Coker}}}
\def\Coim{\mathop{\mathrm{Coim}}}

% Boxtimes
%
\DeclareMathSymbol{\boxtimes}{\mathbin}{AMSa}{"02}

%
% Macros for moduli stacks/spaces
%
\def\QCohstack{\mathcal{QC}\!\mathit{oh}}
\def\Cohstack{\mathcal{C}\!\mathit{oh}}
\def\Spacesstack{\mathcal{S}\!\mathit{paces}}
\def\Quotfunctor{\mathrm{Quot}}
\def\Hilbfunctor{\mathrm{Hilb}}
\def\Curvesstack{\mathcal{C}\!\mathit{urves}}
\def\Polarizedstack{\mathcal{P}\!\mathit{olarized}}
\def\Complexesstack{\mathcal{C}\!\mathit{omplexes}}
% \Pic is the operator that assigns to X its picard group, usage \Pic(X)
% \Picardstack_{X/B} denotes the Picard stack of X over B
% \Picardfunctor_{X/B} denotes the Picard functor of X over B
\def\Pic{\mathop{\mathrm{Pic}}\nolimits}
\def\Picardstack{\mathcal{P}\!\mathit{ic}}
\def\Picardfunctor{\mathrm{Pic}}
\def\Deformationcategory{\mathcal{D}\!\mathit{ef}}

\setlength{\emergencystretch}{3em}
\begin{document}
\title[Stacks Project, Tag 092S]{414 --- Weak dimension one and local valuation-ring structure}
\maketitle
\noindent Copyright (C) 2005--2025 Johan de Jong. Stacks Project authors. Source: \href{https://stacks.math.columbia.edu/tag/092S}{Tag 092S}.

\noindent Editorial difficulty note (not author text). Difficulty H2: The proof passes from flat ideals to arbitrary submodules of flat modules using filtered free approximations and a filtration by ideals. It then uses local finite-flat freeness to prove the local ring is a domain with principal finitely generated ideals, establishing the valuation-ring characterization in both directions..

\noindent Editorial provenance note (not author text). History: excerpt from revision \texttt{a0444\allowbreak 6e57e\allowbreak c1fbc\allowbreak 252a8\allowbreak 71afc\allowbreak ec775\allowbreak 2fb28\allowbreak 07b14}, more-algebra.tex, lines 30574--30645. Reference presentation was adapted; original-author context and core support blocks were retained or added. The mathematical wording is unchanged. Licensed under GNU FDL 1.2 or later; no invariant sections or cover texts. See ../licenses/COPYING. Context blocks reproduce author definitions and supporting results; they are not additional counted problems.

\begin{lemma}
\label{lemma-weak-dimension-at-most-1}
Let $A$ be a ring. The following are equivalent
\begin{enumerate}
\item $A$ has weak dimension $\leq 1$,
\item every ideal of $A$ is flat,
\item every finitely generated ideal of $A$ is flat,
\item every submodule of a flat $A$-module is flat, and
\item every local ring of $A$ is a valuation ring.
\end{enumerate}
\end{lemma}

\begin{proof}
If $A$ has weak dimension $\leq 1$, then the resolution
$0 \to I \to A \to A/I \to 0$ shows that every ideal $I$
is flat by Lemma \href{https://stacks.math.columbia.edu/tag/0653}{lemma last one flat (Tag 0653)}.
Hence (1) $\Rightarrow$ (2).

\medskip\noindent
Assume (4). Let $M$ be an $A$-module. Choose a surjection
$F \to M$ where $F$ is a free $A$-module. Then $\Ker(F \to M)$
is flat by assumption, and we see that $M$ has tor dimension
$\leq 1$ by Lemma \href{https://stacks.math.columbia.edu/tag/066F}{lemma tor dimension (Tag 066F)}.
Hence (4) $\Rightarrow$ (1).

\medskip\noindent
Every ideal is the union of the finitely generated ideals
contained in it. Hence (3) implies (2) by
Algebra, Lemma \href{https://stacks.math.columbia.edu/tag/05UT}{lemma colimit flat (Tag 05UT)}.
Thus (3) $\Leftrightarrow$ (2).

\medskip\noindent
Assume (2). Suppose that $N \subset M$ with $M$ a flat $A$-module.
We will prove that $N$ is flat.
We can write $M = \colim M_i$ with each $M_i$ finite free, see
Algebra, Theorem \href{https://stacks.math.columbia.edu/tag/058G}{theorem lazard (Tag 058G)}.
Setting $N_i \subset M_i$ the inverse image of $N$ we see that
$N = \colim N_i$. By
Algebra, Lemma \href{https://stacks.math.columbia.edu/tag/05UT}{lemma colimit flat (Tag 05UT)}.
it suffices to prove $N_i$ is flat and we reduce
to the case $M = R^{\oplus n}$. In this case
the module $N$ has a finite filtration by the submodules
$R^{\oplus j} \cap N$ whose subquotients are ideals.
By (2) these ideals are flat and hence $N$ is flat by
Algebra, Lemma \href{https://stacks.math.columbia.edu/tag/00HM}{lemma flat ses (Tag 00HM)}. Thus (2) $\Rightarrow$ (4).

\medskip\noindent
Assume $A$ satisfies (1) and let $\mathfrak p \subset A$ be a
prime ideal. By
Lemmas \href{https://stacks.math.columbia.edu/tag/092N}{lemma when weakly etale (Tag 092N)} and \href{https://stacks.math.columbia.edu/tag/092E}{lemma weak dimension goes up (Tag 092E)}
we see that $A_\mathfrak p$ satisfies (1). We will show $A$ is a valuation ring
if $A$ is a local ring satisfying (3). Let $f \in \mathfrak m$
be a nonzero element. Then $(f)$ is a flat nonzero module generated by
one element. Hence it is a free $A$-module by
Algebra, Lemma \href{https://stacks.math.columbia.edu/tag/00NZ}{lemma finite flat local (Tag 00NZ)}.
It follows that $f$ is a nonzerodivisor and $A$ is a domain.
If $I \subset A$ is a finitely generated ideal, then we similarly
see that $I$ is a finite free $A$-module, hence (by considering the
rank) free of rank $1$ and $I$ is a principal ideal. Thus $A$ is a
valuation ring by
Algebra, Lemma \href{https://stacks.math.columbia.edu/tag/090Q}{lemma characterize valuation ring (Tag 090Q)}.
Thus (1) $\Rightarrow$ (5).

\medskip\noindent
Assume (5). Let $I \subset A$ be a finitely generated ideal.
Then $I_\mathfrak p \subset A_\mathfrak p$ is a finitely generated ideal
in a valuation ring, hence principal
(Algebra, Lemma \href{https://stacks.math.columbia.edu/tag/090Q}{lemma characterize valuation ring (Tag 090Q)}), hence flat.
Thus $I$ is flat by
Algebra, Lemma \href{https://stacks.math.columbia.edu/tag/00HT}{lemma flat localization (Tag 00HT)}.
Thus (5) $\Rightarrow$ (3). This finishes the proof of the lemma.
\end{proof}
\end{document}
